% Part: computability
% Chapter: computability-theory
% Section: computable-sets

\documentclass[../../../include/open-logic-section]{subfiles}

\begin{document}

\olfileid{cmp}{thy}{cps}
\olsection{গণনসাধ্য সেট}

গণনসাধ্য অপেক্ষক থেকে গণনসাধ্য সেটে গণনসাধ্যতার ধারণাটি প্রসারিত করা যায়:

\begin{defn}
  ধরা যাক, $S$ স্বাভাবিক সংখ্যার একটি সেট। তাহলে $S$ \emph{গণনসাধ্য}
  যদি এবং কেবল যদি তার চরিত্রসূচক অপেক্ষক~$\Char{S}$ গণনসাধ্য। অন্যভাবে
  বললে, $S$~গণনসাধ্য যদি এবং কেবল যদি অপেক্ষক
\[
\Char{S}(x) =
\begin{cases}
1 & \text{যদি $x \in S$} \\
0 & \text{অন্যথায়}
\end{cases}
\]
গণনসাধ্য। একইভাবে, কোনো সম্বন্ধ $R(x_0, \dots, x_{k-1})$ গণনসাধ্য যদি
এবং কেবল যদি তার চরিত্রসূচক অপেক্ষক গণনসাধ্য।

গণনসাধ্য সেট ও সম্বন্ধকে \emph{নির্ণেয়}-ও বলা হয়।
\end{defn}

\begin{explain}
লক্ষ করুন, এখন আমাদের গণনসাধ্যতার কয়েকটি ধারণা আছে: আংশিক অপেক্ষকের,
অপেক্ষকের এবং সেটের জন্য। এগুলি গুলিয়ে ফেলবেন না!{}
\iftag{TMs}{কোনো আংশিক অপেক্ষক গণনাকারী টুরিং যন্ত্রটি যেসব নিবেশে
  অপেক্ষকটি সংজ্ঞায়িত, সেখানে তার নির্গম ফেরত দেয়; কোনো সেট গণনাকারী
  টুরিং যন্ত্রটি নিবেশটি সেটে আছে কি না নির্ণয় করে~$1$ অথবা~$0$ ফেরত
  দেয়।}{}
\end{explain}

\end{document}
